% ── CHAPTER 15: WHAT IT MEANS (~8,000 words) ──
\chapter{What It Means}
\label{ch:what-it-means}

% STATUS: TODO
% TARGET: ~8,000 words
% DIFFICULTY: Medium — synthesis chapter; needs all other chapters done first.
%   Sources: pub2 §8 (conclusion), pub3 §7 (conclusion), but mostly new writing.
% PRIMARY SOURCES:
%   pub2/sections/08-conclusion.tex  — synthesis of what convergence establishes
%   pub3/sections/07-conclusion.tex  — two-wings confirmation
% NOTE: The final chapter. Addresses four audiences. Ends with the book's
%   closing statement: "Three instruments, one signal."
% HARD RULES: See BOOK_SHELL.md

% \epigraph{The most beautiful thing we can experience is the mysterious.
%   It is the source of all true art and science.}{Albert Einstein}

%  OPENING (~500 words)
% TODO: Fourteen chapters and a paradigm shift later, what has been established?
% State the thesis one final time, completely and precisely:
%   Three independent traditions — quantum physics, process philosophy,
%   Bahá'í metaphysics — have arrived at the same structural description of
%   reality: that existence is constituted by relationship, and that the force
%   sustaining relational coherence is, in each tradition's vocabulary, love.
%   The convergence is not analogical but structural. Not borrowed but independent.
%   Not approximate but precise enough to formalize.

\section{What It Means for Science}
\label{sec:ch15-science}
% (~1,500 words)
% TODO: The implications for scientific research programs.
% KEY POINTS:
%   - The RCF provides a new explanatory/operational framework for decoherence
%     and entanglement in quantum information science
%   - The convergence provides a new way of thinking about the interpretive
%     problem in QM: the relational interpretation is not just "one option
%     among many" — it has two independent traditions pointing in the same direction
%   - The paradigm shift (ch13a) points toward research directions:
%     information-first ontology, non-composite entities, dual-aspect frameworks
%   - The decoherence objection against the non-composite soul is a category
%     error — which means there is genuine philosophical work to be done on
%     what our physical theories miss
%   - Science and religion are not in conflict here; they are in convergence
% SOURCE: pub2 §8; write fresh
% CITE: AdamsRCT, AdamsThreePaths

\section{What It Means for Philosophy}
\label{sec:ch15-philosophy}
% (~1,500 words)
% TODO: The implications for philosophy of mind and metaphysics.
% KEY POINTS:
%   - The mind-body problem looks different from a relational ontology:
%     not "how does physical stuff produce mental stuff" but "what is the
%     relational structure of consciousness?"
%   - The decoherence argument against the soul is a category error, not
%     a refutation — which reopens the question of immaterial substance
%   - Process philosophy receives external confirmation from two independent
%     traditions — which matters for its reception as a metaphysical framework
%   - The hard problem of consciousness: the transceiver model is one way
%     of approaching it, not as elimination but as reframing
% POSSIBLE CITATION: Chalmers on hard problem (TODO: add bib entry if used)
% SOURCE: pub2 §8; write fresh

\section{What It Means for the Science-Religion Question}
\label{sec:ch15-science-religion}
% (~1,500 words)
% TODO: The implications for the science-religion relationship.
% KEY POINTS:
%   - The science-religion conflict was always a conflict between SPECIFIC
%     interpretations of specific findings, not between science and religion
%     as such
%   - What this book documents is what genuine non-conflict looks like:
%     two independent methods (empirical science and revelation) detecting
%     the same structural feature of reality
%   - This is what the two-wings principle predicts — and what three centuries
%     of assuming they must conflict obscured
%   - The Bahá'í principle of the harmony of religion and science is not
%     wishful thinking or diplomatic compromise — it is an empirical prediction
%     that this book has, in one domain, verified
% SOURCE: pub3 §7 (conclusion); write fresh
% CITE: AdamsMahabbat, PUP, ParisTalks

\section{What It Means for You}
\label{sec:ch15-you}
% (~1,500 words)
% TODO: The most personal section of the book. Address the individual reader.
% KEY POINTS:
%   - Love is not a sentiment about something else. It is the name for the
%     most fundamental force in reality, experienced at its most conscious
%     and intense level in human relationships.
%   - When you love — with full attention, full presence, full relational
%     coherence — you are participating in the same principle that holds
%     quantum systems in coherence across space and constitutes existence at
%     every degree.
%   - This is not a mystical claim dressed in scientific language. It is
%     the conclusion of three independent traditions arriving at the same
%     structural description of reality.
%   - The question "does love matter?" has a precise answer: yes, because
%     love IS what matters — what constitutes the relational fabric of reality.
%   - Live accordingly.
% SOURCE: Write fresh; draw on pub2 §6 final paragraphs and ch14 §4
% NOTE: This section should be beautiful. It is the reason the book was written.

%  CLOSING (~500 words)
% TODO: The closing statement. "Three instruments, one signal."
% Return to the metaphor of the convergence: each tradition is an instrument
% pointed at reality. When three independent instruments — each with its own
% methodology, vocabulary, and calibration — register the same signal, the
% signal is real.
% The Physics of Love is not a poetic title. It is a claim:
%   that the force that holds atoms together, that constitutes existence
%   at every degree, that quantum physics measures with the Affinity Tensor,
%   that Whitehead calls eros, that the Bahá'í writings call mahabbat —
%   is the same force that, in its most conscious and fully experienced form,
%   human beings call love.
% Three instruments. One signal.

\begin{convergencemoment}{What It Means}
% TODO: The final Convergence Moment — the whole book in one box.
% From the quantum side: Bell's theorem closes local substance ontology;
%   relational constitution of properties; Affinity Tensor as formal
%   representation of the coherence that constitutes existence.
% From the process-philosophical side: actual occasions are constituted
%   by prehension; eros is the creative drive toward richer relational
%   integration; process is primary; substance is derivative.
% From the Bahá'í side: existence = relational affinity; motion = life;
%   the creative foundation at every degree is an expression of love.
%   The soul is non-composite and relationally embedded.
% Three vocabularies, one signal: the universe is, at its most fundamental
%   level, a relational structure held together by love.
\end{convergencemoment}
